\documentclass[11pt]{article}
\usepackage{amsmath}
\usepackage[a5paper,margin=18mm]{geometry}

\title{Bilan de masse d'un glacier alpin}
\author{Atelier}
\date{}

\begin{document}
\maketitle

\section{Introduction}
Les glaciers de montagne perdent de la masse depuis plusieurs
d\'ecennies. Cette note sert de document d'essai pour la
synchronisation entre la source et le PDF compil\'e.

Le bilan de masse annuel $B_a$ s'\'ecrit comme la somme de
l'accumulation hivernale et de l'ablation estivale, soit
$B_a = B_w + B_s$ avec $B_s < 0$ la plupart des ann\'ees.

% Une ligne de commentaire qui ne produit rien.

\begin{equation}
  B_a = \int_{t_0}^{t_1} \dot{b}(t)\,\mathrm{d}t
  \label{eq:bilan}
\end{equation}

\section{Donn\'ees}
\begin{itemize}
  \item Mesures de balises d'ablation, relev\'ees chaque automne.
  \item Sondages de neige au printemps, sur un r\'eseau fixe.
  \item Alb\'edo MODIS journalier, filtr\'e par la couverture nuageuse.
\end{itemize}

\section{M\'ethode}
La m\'ethode glaciologique combine les relev\'es ponctuels et une
interpolation par bandes d'altitude. Chaque bande re\c{c}oit le
bilan moyen de ses balises, pond\'er\'e par sa surface.

\subsection{Interpolation}
Pour une bande $i$ de surface $S_i$ et de bilan $b_i$, le bilan
sp\'ecifique du glacier vaut
\[
  B = \frac{\sum_i S_i\, b_i}{\sum_i S_i}.
\]
Cette moyenne pond\'er\'ee suppose que le bilan varie peu \`a
l'int\'erieur d'une bande, hypoth\`ese raisonnable pour des bandes
de cent m\`etres.

\begin{table}[h]
  \centering
  \begin{tabular}{lrr}
    Bande & Surface (km$^2$) & Bilan (m w.e.) \\
    2800--2900 & 0.42 & $-2.1$ \\
    2900--3000 & 0.77 & $-1.4$ \\
    3000--3100 & 0.91 & $-0.6$ \\
  \end{tabular}
  \caption{Bilan par bande d'altitude.}
\end{table}

\subsection{Incertitudes}
Les erreurs de mesure des balises sont de l'ordre de dix
centim\`etres d'\'equivalent en eau\footnote{Valeur typique pour
des balises en bois relev\'ees au ruban.}. L'erreur
d'interpolation domine largement pour les bandes sans balise.

\newpage
\subsection{Validation}
Une validation g\'eod\'esique, par diff\'erence de mod\`eles
num\'eriques de terrain, compl\`ete la m\'ethode glaciologique
sur des p\'eriodes de cinq \`a dix ans.

\begin{figure}[h]
  \centering
  \rule{30mm}{12mm}\hspace{4mm}%
  \fbox{\parbox{36mm}{Balise B7, relev\'ee
    en septembre, fond\'ee dans la glace.}}
  \caption{Sch\'ema d'une balise.}
\end{figure}

\noindent\begin{minipage}[t]{0.45\linewidth}
  Colonne de gauche : les mesures
  de l'ann\'ee hydrologique.
\end{minipage}\hfill
\begin{minipage}[t]{0.45\linewidth}
  Colonne de droite : la comparaison
  avec la s\'erie longue.
\end{minipage}


\section{Conclusion}
Le bilan reste n\'egatif sur toute la p\'eriode, ce qui confirme
la tendance r\'egionale d\'ej\`a d\'ecrite dans la litt\'erature.
Les incertitudes principales viennent de la densit\'e de la neige
et de la repr\'esentativit\'e spatiale des balises.

Un dernier paragraphe, assez long pour occuper plusieurs lignes de
texte dans la colonne \'etroite du format A5, permet de v\'erifier
qu'un clic au milieu d'un paragraphe renvoie bien vers la bonne
ligne source, et non vers le d\'ebut ou la fin du paragraphe entier.

\end{document}
